\chapter{Grafos}


\section{Representaciones comunes}

\subsection*{Lista de adyacencia}

\begin{lstlisting}[caption=Lista de adyacencia]
	vector<int> G[MAXN]; // grafo sin peso
	
	void agregar_arista(int u, int v) {
		G[u].push_back(v);
		G[v].push_back(u); // si es no dirigido
	}
\end{lstlisting}

\subsection*{Matriz de adyacencia}

\begin{lstlisting}[caption=Matriz de adyacencia]
	bool mat[MAXN][MAXN];
	
	void agregar_arista(int u, int v) {
		mat[u][v] = true;
		mat[v][u] = true; // si es no dirigido
	}
\end{lstlisting}

\section{Recorridos clásicos}

\subsection{DFS (Depth-First Search)}

\begin{lstlisting}[caption=DFS no dirigido]
	bool vis[MAXN];
	vector<int> G[MAXN];
	
	void dfs(int u) {
		vis[u] = true;
		for (int v : G[u]) {
			if (!vis[v]) dfs(v);
		}
	}
\end{lstlisting}

\subsection{BFS (Breadth-First Search)}

\begin{lstlisting}[caption=BFS con cálculo de distancias]
	int dist[MAXN];
	bool vis[MAXN];
	vector<int> G[MAXN];
	
	void bfs(int s) {
		queue<int> q;
		q.push(s);
		vis[s] = true;
		dist[s] = 0;
		
		while (!q.empty()) {
			int u = q.front(); q.pop();
			for (int v : G[u]) {
				if (!vis[v]) {
					vis[v] = true;
					dist[v] = dist[u] + 1;
					q.push(v);
				}
			}
		}
	}
\end{lstlisting}

\section{Componentes conexas}

\begin{lstlisting}[caption=Contar componentes conexas]
	int comp = 0;
	
	for (int i = 0; i < n; ++i) {
		if (!vis[i]) {
			dfs(i);
			comp++;
		}
	}
\end{lstlisting}

\section{Dijkstra (caminos mínimos)}

\begin{lstlisting}[caption=Dijkstra con priority queue]
	typedef pair<int, int> pii;
	vector<pii> G[MAXN]; // {peso, destino}
	int dist[MAXN];
	
	void dijkstra(int s) {
		priority_queue<pii, vector<pii>, greater<pii>> pq;
		fill(dist, dist + MAXN, INF);
		dist[s] = 0;
		pq.push({0, s});
		
		while (!pq.empty()) {
			auto [d, u] = pq.top(); pq.pop();
			if (d > dist[u]) continue;
			for (auto [w, v] : G[u]) {
				if (dist[u] + w < dist[v]) {
					dist[v] = dist[u] + w;
					pq.push({dist[v], v});
				}
			}
		}
	}
\end{lstlisting}

\section{Bellman-Ford}

\begin{lstlisting}[caption=Bellman-Ford para grafos con pesos negativos]
	vector<tuple<int, int, int>> edges; // {u, v, peso}
	int dist[MAXN];
	
	bool bellman(int s, int n) {
		fill(dist, dist + MAXN, INF);
		dist[s] = 0;
		
		for (int i = 0; i < n - 1; ++i) {
			for (auto [u, v, w] : edges) {
				if (dist[u] + w < dist[v]) {
					dist[v] = dist[u] + w;
				}
			}
		}
		
		// Detectar ciclos negativos
		for (auto [u, v, w] : edges) {
			if (dist[u] + w < dist[v]) return true;
		}
		
		return false;
	}
\end{lstlisting}

\section{Problemas clásicos}

\subsection*{A - Número de componentes conexas}

\textbf{Estado:} nodos visitados.  
\textbf{Recorrido:} DFS o BFS.  
\textbf{Complejidad:} \( O(n + m) \)

\subsection*{B - Camino más corto (no ponderado)}

\textbf{Idea:} usar BFS desde el nodo origen y guardar las distancias.

\subsection*{C - Dijkstra / Bellman-Ford}

\textbf{Cuando usar cada uno:}

- Dijkstra → aristas no negativas.
- Bellman-Ford → permite pesos negativos.

\section{¿Cómo verificar si un grafo es una línea (path)?}



\begin{enumerate}
	\item Verificar que \( M = N - 1 \) (condición necesaria para ser un árbol).
	\item Verificar que el grafo es conexo (DFS o BFS desde cualquier nodo).
	\item Contar los grados de los nodos:
	\begin{itemize}
		\item Debe haber exactamente 2 nodos con grado 1.
		\item Todos los demás deben tener grado 2.
	\end{itemize}
\end{enumerate}

\subsection*{Código en C++}

\begin{lstlisting}[caption=Verificar si un grafo es una línea]
	vector<int> G[MAXN];
	bool vis[MAXN];
	int deg[MAXN];
	
	void dfs(int u) {
		vis[u] = true;
		for (int v : G[u]) {
			if (!vis[v]) dfs(v);
		}
	}
	
	bool es_linea(int n, int m) {
		if (m != n - 1) return false; // debe tener N-1 aristas
		
		// verificar conectividad
		fill(vis, vis + n, false);
		dfs(0);
		for (int i = 0; i < n; ++i)
		if (!vis[i]) return false;
		
		// contar grados
		int cnt_deg1 = 0;
		for (int i = 0; i < n; ++i) {
			if (deg[i] == 1) cnt_deg1++;
			else if (deg[i] != 2) return false;
		}
		
		return cnt_deg1 == 2;
	}
\end{lstlisting}

\textbf{Entrada:}
\begin{itemize}
	\item Llenar \texttt{G[i]} con las aristas del grafo.
	\item Aumentar \texttt{deg[u]} y \texttt{deg[v]} cada vez que se agrega una arista entre \( u \) y \( v \).
\end{itemize}
		
\section{Detectar ciclos en grafo no dirigido}

\begin{enumerate}
	\item Usar DFS manteniendo el nodo padre
	\item Si encuentras un nodo ya visitado que no es el padre, hay ciclo
\end{enumerate}

\subsection*{Código en C++}

\begin{lstlisting}[caption=Detectar ciclos en grafo no dirigido]
	vector<int> G[MAXN];
	bool vis[MAXN];
	
	bool hasCycle(int u, int parent) {
		vis[u] = true;
		for (int v : G[u]) {
			if (!vis[v]) {
				if (hasCycle(v, u)) 
					return true;
			}
			else if (v != parent) 
				return true;
		}
		return false;
	}
	
	bool hasCycleAll(int n) {
		memset(vis, 0, sizeof(vis));
		for (int i = 0; i < n; ++i) {
			if (!vis[i]) {
				if (hasCycle(i, -1)) {
					return true;
				}
			}
		}
		return false;
	}
\end{lstlisting}

\section{Detectar ciclos en grafo dirigido}

Inicialmente, todos los nodos están en blanco (0), lo que significa que no fueron visitados. Cuando empezamos a explorar un nodo, lo marcamos como gris (1), indicando que está en proceso dentro de la recursión. Luego recorremos sus vecinos: si encontramos uno blanco, seguimos el DFS; pero si encontramos un nodo gris, significa que volvimos a un nodo que todavía está en la pila de recursión, lo que implica la existencia de un ciclo (una arista de retroceso o back edge). Finalmente, cuando terminamos de procesar un nodo, lo marcamos como negro (2), indicando que ya fue completamente explorado. Para asegurarnos de detectar ciclos en todo el grafo, este proceso se inicia desde cada nodo no visitado. La función \texttt{hasCycleAll()} es la que se usaría en el \texttt{main} para verificar que tu grafo tiene un ciclo.

\begin{enumerate}
	\item Usar DFS con colores (0: no visitado, 1: en proceso, 2: visitado)
	\item Si encuentras un nodo "en proceso" durante el DFS, hay ciclo
\end{enumerate}

\subsection*{Código en C++}

\begin{lstlisting}[caption=Detectar ciclos en grafo dirigido]
	vector<int> G[MAXN];
	// 0: blanco (no visitado)
	// 1: gris (en proceso)
	// 2: negro (visitado)
	int color[MAXN]; 
	
	bool hasCycleDirected(int u) {
		color[u] = 1; // gris
		for (int v : G[u]) {
			if (color[v] == 0) {
				if (hasCycleDirected(v)) 
					return true;
			} 
			else if (color[v] == 1) 
				return true;
		}
		color[u] = 2; // negro
		return false;
	}
	
	bool hasCycleAll(int n) {
		fill(color, color + n, 0);
		for (int i = 0; i < n; ++i) {
			if (color[i] == 0 && hasCycleDirected(i))
				return true;
		}
		return false;
	}
\end{lstlisting}

\section{Verificar si hay ciclo impar}

\begin{enumerate}
	\item Intentar colorear el grafo con 2 colores (bipartición)
	\item Si no es posible, contiene ciclo impar
\end{enumerate}

\subsection*{Código en C++}

\begin{lstlisting}[caption=Detectar ciclo impar]
	vector<int> G[MAXN];
	int col[MAXN]; // 0 o 1
	
	bool tieneCicloImpar(int n) {
		queue<int> q;
		fill(col, col + n, -1);
		
		col[0] = 0;
		q.push(0);
		
		while (!q.empty()) {
			int u = q.front(); q.pop();
			for (int v : G[u]) {
				if (col[v] == -1) {
					col[v] = 1 - col[u];
					q.push(v);
				}
				else if (col[v] == col[u]) {
					return true; // Ciclo impar encontrado
				}
			}
		}
		return false;
	}
\end{lstlisting}

\section{Verificar si el grafo es bipartito}

\begin{enumerate}
	\item Igual que detectar ciclo impar, pero devolviendo el resultado inverso
	\item Si se puede colorear con 2 colores, es bipartito
\end{enumerate}

\subsection*{Código en C++}

\begin{lstlisting}[caption=Verificar si grafo es bipartito]
	vector<int> G[MAXN];
	int col[MAXN]; // 0 o 1
	
	bool esBipartito(int n) {
		queue<int> q;
		fill(col, col + n, -1);
		
		col[0] = 0;
		q.push(0);
		
		while (!q.empty()) {
			int u = q.front(); q.pop();
			for (int v : G[u]) {
				if (col[v] == -1) {
					col[v] = 1 - col[u];
					q.push(v);
				}
				else if (col[v] == col[u]) {
					return false;
				}
			}
		}
		return true;
	}
\end{lstlisting}

\section{Distancia mínima desde un nodo}

\begin{enumerate}
	\item Usar BFS para calcular distancias desde el nodo fuente
	\item Mantener arreglo de distancias inicializado en -1
\end{enumerate}

\subsection*{Código en C++}

\begin{lstlisting}[caption=Distancia mínima con BFS]
	vector<int> G[MAXN];
	int dist[MAXN];
	
	void bfs(int s, int n) {
		queue<int> q;
		fill(dist, dist + n, -1);
		
		dist[s] = 0;
		q.push(s);
		
		while (!q.empty()) {
			int u = q.front(); q.pop();
			for (int v : G[u]) {
				if (dist[v] == -1) {
					dist[v] = dist[u] + 1;
					q.push(v);
				}
			}
		}
	}
\end{lstlisting}

\section{Verificar si el grafo es un árbol}

\begin{enumerate}
	\item Verificar que tiene exactamente N-1 aristas
	\item Verificar que es conexo
	\item Verificar que no tiene ciclos
\end{enumerate}

\subsection*{Código en C++}

\begin{lstlisting}[caption=Verificar si grafo es árbol]
	vector<int> G[MAXN];
	bool vis[MAXN];
	
	bool hasCycle(int u, int parent) {
		vis[u] = true;
		for (int v : G[u]) {
			if (!vis[v]) {
				if (hasCycle(v, u)) return true;
			}
			else if (v != parent) return true;
		}
		return false;
	}
	
	bool esArbol(int n, int m) {
		if (m != n - 1) return false;
		
		// Verificar conexidad
		fill(vis, vis + n, false);
		hasCycle(0, -1);
		for (int i = 0; i < n; ++i)
		if (!vis[i]) return false;
		
		// Verificar que no tiene ciclos
		fill(vis, vis + n, false);
		if (hasCycle(0, -1)) return false;
		
		return true;
	}
\end{lstlisting}

\section{Número de hojas en un árbol}

\begin{enumerate}
	\item Contar nodos con grado 1 (excepto posiblemente la raíz en árboles enraizados)
\end{enumerate}

\subsection*{Código en C++}

\begin{lstlisting}[caption=Contar hojas en un árbol]
	vector<int> G[MAXN];
	int deg[MAXN];
	
	int contarHojas(int n) {
		int hojas = 0;
		for (int i = 0; i < n; ++i) {
			if (G[i].size() == 1) // grado 1
			hojas++;
		}
		return hojas;
	}
\end{lstlisting}

\section{Reconstrucción de Caminos (BFS)}

Cuando necesitas no solo la distancia mínima, sino la secuencia de nodos, se utiliza un arreglo de predecesores (o padres).

\subsection*{1. Concepto Clave: El arreglo parent}

\begin{itemize}
	\item \texttt{parent[v] = u}: Significa que llegamos al nodo \texttt{v} desde el nodo \texttt{u}.
	\item Inicializamos todo el arreglo en -1 (o un valor nulo) para identificar el origen y los nodos no visitados.
\end{itemize}

\subsection*{2. Implementación del BFS}

\begin{lstlisting}[caption=BFS con arreglo parent]
// n = numero de nodos, start = nodo inicial
vector<int> parent(n + 1, -1);
vector<int> dist(n + 1, -1);
queue<int> q;

q.push(start);
dist[start] = 0;

while(!q.empty()){
	int u = q.front(); q.pop();
	
	for(int v : G[u]){
		if(dist[v] == -1){ // Si no ha sido visitado
			dist[v] = dist[u] + 1;
			parent[v] = u; // REGISTRO: u es el padre de v
			q.push(v);
		}
	}
}
\end{lstlisting}

\subsection*{3. Función para recuperar el camino}

Una vez termina el BFS, el camino está encadenado desde el destino hacia atrás. Debes seguir los punteros hasta llegar al origen y luego invertir la lista.

\begin{lstlisting}[caption=Función para reconstruir el camino]
vector<int> getPath(int start, int end, const vector<int>& parent) {
	vector<int> path;
	
	// Si el destino no tiene padre y no es el inicio, no hay camino
	if (parent[end] == -1 && start != end) return {};

	// Retrocedemos desde 'end' hasta 'start'
	for (int v = end; v != -1; v = parent[v]) {
		path.push_back(v);
	}
	
	// El camino esta [destino, ..., origen], lo invertimos
	reverse(path.begin(), path.end());
	
	return path;
}
\end{lstlisting}

\section{Orden Topológico}

Se te proporciona un grafo dirigido con $n$ vértices y $m$ aristas. Tienes que encontrar un orden de los vértices, de modo que cada arista vaya desde el vértice con un índice menor hasta un vértice con un índice mayor.

En otras palabras, se trata de encontrar una permutación de los vértices (orden topológico) que corresponda al orden definido por todas las aristas del grafo. Un orden topológico puede no existir en absoluto. Solo existe si el grafo dirigido no contiene ciclos (DAG).

\subsection*{Ejemplo de uso}

Un problema común en el que se produce la ordenación topológica es el siguiente. Hay $n$ variables con valores desconocidos. Para algunas variables, sabemos que una de ellas es menor que la otra. Debes comprobar si estas restricciones son contradictorias y, si no lo son, mostrar las variables en orden ascendente (si hay varias respuestas posibles, mostrar cualquiera de ellas). Es fácil notar que este es exactamente el problema de encontrar el orden topológico de un grafo con $n$ vértices.

\subsection*{¿Cómo funciona con DFS?}

Supongamos que el grafo es acíclico. ¿Qué hace la búsqueda en profundidad (DFS)?

Al comenzar desde algún vértice $v$, DFS intenta recorrer a lo largo de todas las aristas que salen de $v$. Se detiene en los bordes cuyos extremos ya han sido visitados previamente, y recorre el resto de los bordes y continúa recursivamente en sus extremos.

Por lo tanto, en el momento en que la llamada a la función $\text{dfs}(v)$ ha terminado, todos los vértices que son alcanzables desde $v$ han sido visitados por la búsqueda, ya sea directamente (a través de una arista) o indirectamente.

Añadamos el vértice $v$ a una lista cuando terminemos $\text{dfs}(v)$. Como todos los vértices alcanzables ya han sido visitados, ya estarán en la lista cuando agreguemos $v$. Hagamos esto para cada vértice del grafo, con una o varias búsquedas en profundidad. Para cada arista dirigida $v \rightarrow u$ en el gráfico, $u$ aparecerá antes en esta lista que $v$, porque $u$ es accesible desde $v$. 

Entonces, si simplemente etiquetamos los vértices en esta lista con $n-1, n-2, \dots, 1, 0$, hemos encontrado un orden topológico del grafo. En otras palabras, la lista representa el orden topológico inverso, por lo que basta con invertir el arreglo al final.

\subsection*{Código en C++}

Aquí se presenta una implementación que asume que el grafo es acíclico, es decir, que existe el orden topológico deseado:

\begin{lstlisting}[caption=Orden Topológico con DFS]
// n = numero de vertices
int n; 
vector<vector<int>> adj; 
vector<bool> visited;
vector<int> ans;

void dfs(int v) {
        visited[v] = true;
        for (int u : adj[v]) {
                if (!visited[u]) {
                        dfs(u);
                }
        }
        ans.push_back(v);
}

void topological_sort() {
        visited.assign(n, false);
        ans.clear();
        for (int i = 0; i < n; ++i) {
                if (!visited[i]) {
                        dfs(i);
                }
        }
        reverse(ans.begin(), ans.end());
}
\end{lstlisting}

La función principal de la solución es \texttt{topological\_sort}, que inicializa las variables DFS, lanza DFS y recibe la respuesta en el vector \texttt{ans}. 

Vale la pena señalar que cuando el grafo \textbf{no} es acíclico, el resultado de \texttt{topological\_sort} aún sería algo significativo en el sentido de que si un vértice $u$ es alcanzable desde el vértice $v$, pero no al revés, el vértice $v$ siempre aparecerá primero en el array resultante.

